\documentclass[12pt,a4paper,oneside]{book}

\usepackage[T1]{fontenc}
\usepackage[utf8]{inputenc}
\usepackage[italian]{babel}
\usepackage{lmodern}
\usepackage{microtype}
\usepackage{amsmath,amssymb,mathtools}
\usepackage{booktabs}
\usepackage{array}
\usepackage{tabularx}
\usepackage{csquotes}
\usepackage[hidelinks]{hyperref}
\usepackage[backend=biber,style=authoryear]{biblatex}

\addbibresource{references.bib}

\title{Introduzione ai paper sulla tail dependence}
\author{Massimo Gatti}
\date{}

\begin{document}

\maketitle

\chapter[Quadro introduttivo]{Tail dependence: quadro introduttivo della letteratura}
\label{chap:quadro-letteratura}

\section{Obiettivo e percorso di lettura}

La dipendenza tra eventi estremi costituisce un problema centrale nella
misurazione del rischio finanziario. Due attività possono mostrare una
correlazione moderata durante le fasi ordinarie del mercato e, allo stesso
tempo, subire perdite molto elevate congiuntamente durante una crisi. Una
misura globale di associazione, come la correlazione lineare, non è quindi
sufficiente a descrivere il comportamento congiunto nelle regioni estreme
della distribuzione.

I tre lavori analizzati in questo capitolo affrontano il problema a livelli
progressivamente più avanzati. Il capitolo di \textcite{schmidt2005} introduce
il coefficiente classico di tail dependence, ne presenta il calcolo per alcune
importanti famiglie distributive e discute la stima e le applicazioni al
Value-at-Risk. \textcite{duranteEtAl2015} studiano il legame tra la sezione
diagonale di una copula e la dipendenza nelle code, proponendo la tail
concentration function per descrivere anche il comportamento a soglie finite.
Infine, \textcite{huaJoe2011} introducono il tail order e le tail order
functions, che permettono di distinguere differenti gradi di dipendenza
residua anche in presenza di indipendenza asintotica.

Questi contributi non rappresentano tre direzioni isolate. Essi delineano un
percorso unitario:

\begin{enumerate}
    \item definizione e stima del coefficiente classico di tail dependence;
    \item analisi dei limiti informativi di un singolo coefficiente asintotico;
    \item introduzione di strumenti capaci di descrivere il comportamento
    pre-asintotico e forme intermedie di dipendenza nelle code.
\end{enumerate}

\section{Copule e separazione tra margini e dipendenza}

Sia $\mathbf{X}=(X_1,\ldots,X_d)$ un vettore aleatorio con funzione di
ripartizione congiunta $F$ e funzioni di ripartizione marginali
$F_1,\ldots,F_d$. Il teorema di Sklar consente di scrivere

\begin{equation}
    F(x_1,\ldots,x_d)
    = C\!\left(F_1(x_1),\ldots,F_d(x_d)\right),
    \label{eq:sklar}
\end{equation}

dove $C\colon[0,1]^d\to[0,1]$ è una copula. Se le distribuzioni marginali
sono continue, la copula è unica. Tale rappresentazione separa le
caratteristiche delle singole variabili dalla loro struttura di dipendenza.
La dipendenza di coda può pertanto essere studiata direttamente attraverso
$C$, indipendentemente dalle unità di misura e dalle trasformazioni
strettamente crescenti applicate ai margini.

Nelle applicazioni finanziarie è necessario chiarire la convenzione adottata.
Se $X_1$ e $X_2$ rappresentano rendimenti, le perdite più gravi si collocano
nella coda inferiore. Se, invece, le variabili rappresentano direttamente
perdite monetarie, gli eventi avversi si collocano nella coda superiore. Le
espressioni ``coda inferiore'' e ``coda superiore'' hanno dunque un significato
probabilistico preciso, ma la loro interpretazione economica dipende dalla
definizione delle variabili.

\section{Il coefficiente classico di tail dependence}

\subsection{Definizione}

Il punto di partenza più accessibile è fornito da \textcite{schmidt2005}. Per
una coppia di variabili continue $(X_1,X_2)$, il coefficiente di dipendenza
nella coda superiore è definito, quando il limite esiste, come

\begin{equation}
    \lambda_U
    = \lim_{u\uparrow 1}
    \Pr\!\left[
        X_1>F_1^{-1}(u)
        \,\middle|\,
        X_2>F_2^{-1}(u)
    \right].
    \label{eq:lambda-upper-def}
\end{equation}

Analogamente, il coefficiente di dipendenza nella coda inferiore è

\begin{equation}
    \lambda_L
    = \lim_{u\downarrow 0}
    \Pr\!\left[
        X_1\leq F_1^{-1}(u)
        \,\middle|\,
        X_2\leq F_2^{-1}(u)
    \right].
    \label{eq:lambda-lower-def}
\end{equation}

I coefficienti appartengono all'intervallo $[0,1]$. Un valore positivo indica
dipendenza asintotica nella coda considerata; un valore nullo indica
indipendenza asintotica. Questa terminologia riguarda il limite per soglie
sempre più estreme e non implica necessariamente indipendenza a livelli
osservabili in un campione finito.

In termini della copula $C$, le definizioni diventano

\begin{align}
    \lambda_L
    &= \lim_{u\downarrow 0}\frac{C(u,u)}{u},
    \label{eq:lambda-lower-copula}\\
    \lambda_U
    &= \lim_{u\uparrow 1}
       \frac{1-2u+C(u,u)}{1-u}.
    \label{eq:lambda-upper-copula}
\end{align}

Le equazioni \eqref{eq:lambda-lower-copula} e
\eqref{eq:lambda-upper-copula} mostrano che i coefficienti dipendono dal
comportamento della copula lungo la diagonale principale in prossimità degli
angoli $(0,0)$ e $(1,1)$.

\subsection{Perché la correlazione non è sufficiente}

Il confronto tra distribuzione normale e distribuzione $t$ è particolarmente
istruttivo. Una copula gaussiana non degenere presenta
$\lambda_L=\lambda_U=0$, indipendentemente dal fatto che la correlazione sia
elevata. Una copula $t$, invece, può presentare dipendenza asintotica positiva
in entrambe le code. Per una copula $t$ bivariata con correlazione $\rho$ e
$\nu$ gradi di libertà,

\begin{equation}
    \lambda_L=\lambda_U
    =2\,T_{\nu+1}\!\left(
        -\sqrt{\frac{(\nu+1)(1-\rho)}{1+\rho}}
    \right),
    \label{eq:t-copula-tail-dependence}
\end{equation}

dove $T_{\nu+1}$ denota la funzione di ripartizione di una variabile $t$ di
Student con $\nu+1$ gradi di libertà. Due modelli con lo stesso parametro di
correlazione possono quindi implicare rischi congiunti estremi molto diversi.

\subsection{Stima e applicazione finanziaria}

Oltre alle definizioni, \textcite{schmidt2005} confronta stimatori parametrici
e non parametrici dei coefficienti di coda. Gli stimatori non parametrici si
basano sulla copula empirica e richiedono la scelta di una soglia: utilizzare
troppo poche osservazioni estreme aumenta la varianza, mentre includerne troppe
introduce distorsione perché ci si allontana dalla regione di coda. Il capitolo
mostra inoltre, mediante dati finanziari e una simulazione, che la struttura di
tail dependence può incidere materialmente sulla stima del Value-at-Risk di un
portafoglio.

\section{Diagonale della copula e dipendenza a scala finita}

\subsection{La sezione diagonale}

Il contributo di \textcite{duranteEtAl2015} parte dalla sezione diagonale

\begin{equation}
    \delta_C(t)=C(t,t), \qquad t\in[0,1].
    \label{eq:diagonal-section}
\end{equation}

Le equazioni \eqref{eq:lambda-lower-copula} e
\eqref{eq:lambda-upper-copula} dipendono esclusivamente da
$\delta_C$. Conoscere la diagonale non identifica necessariamente tutta la
copula: copule differenti possono condividere la stessa sezione diagonale.
Essa contiene tuttavia informazione essenziale sul verificarsi congiunto di
eventi che occupano lo stesso livello quantile nei due margini.

Gli autori riesaminano le condizioni che una funzione deve soddisfare per
essere la diagonale di una copula e discutono costruzioni di copule aventi una
diagonale prescritta. Questa parte chiarisce una distinzione importante: la
conoscenza del comportamento di coda lungo la diagonale vincola la struttura
di dipendenza, ma in generale non determina univocamente il modello
multivariato.

\subsection{La tail concentration function}

Un singolo limite può risultare poco informativo nei campioni reali, perché i
dati non raggiungono mai il punto asintotico. Per descrivere la dipendenza a
livelli finiti, \textcite{duranteEtAl2015} considerano la tail concentration
function

\begin{equation}
    q_C(t)=
    \begin{cases}
        \displaystyle \frac{C(t,t)}{t},
        & 0<t\leq \tfrac{1}{2},\\[8pt]
        \displaystyle \frac{1-2t+C(t,t)}{1-t},
        & \tfrac{1}{2}<t<1.
    \end{cases}
    \label{eq:tail-concentration-function}
\end{equation}

Alle estremità del dominio, quando i limiti esistono,

\begin{equation}
    \lim_{t\downarrow 0}q_C(t)=\lambda_L,
    \qquad
    \lim_{t\uparrow 1}q_C(t)=\lambda_U.
\end{equation}

La funzione $q_C$ conserva quindi i coefficienti classici, ma mostra anche il
percorso di convergenza verso di essi. Due copule con
$\lambda_L=\lambda_U=0$ possono avere tail concentration functions molto
diverse e generare probabilità congiunte non trascurabili alle soglie
effettivamente utilizzate nel risk management.

Data una coppia di osservazioni $(X_i,Y_i)_{i=1}^n$, la copula empirica
$\widehat C_n$ consente di stimare la funzione sostituendo $C$ con
$\widehat C_n$ nell'equazione \eqref{eq:tail-concentration-function}. Gli
autori impiegano la funzione stimata per costruire misure di dissimilarità tra
copule, rappresentazioni grafiche dello spazio dei modelli e procedure di
clustering per serie finanziarie. Il principale messaggio applicativo è che
la valutazione del rischio dovrebbe considerare congiuntamente comportamento
asintotico e comportamento a scala finita.

\section{Tail order e dipendenza intermedia}

\subsection{Il limite della classificazione binaria}

La distinzione $\lambda>0$ contro $\lambda=0$ separa dipendenza e indipendenza
asintotica, ma tratta come equivalenti molti modelli che equivalenti non sono.
Se due copule hanno coefficiente nullo, le rispettive probabilità congiunte di
coda possono comunque convergere a zero a velocità differenti. Il contributo
di \textcite{huaJoe2011} formalizza questa informazione mediante il tail order.

Per la coda inferiore di una copula $d$-dimensionale, gli autori considerano
una relazione asintotica della forma

\begin{equation}
    C(u,\ldots,u)
    \sim u^{\kappa_L}\ell_L(u),
    \qquad u\downarrow 0,
    \label{eq:lower-tail-order}
\end{equation}

dove $\ell_L$ è lentamente variabile in zero. Una definizione analoga per la
coda superiore è formulata mediante la survival copula. L'esponente
$\kappa_L$ determina la velocità dominante con cui la probabilità congiunta
si annulla.

Sotto opportune condizioni di dipendenza positiva,

\begin{equation}
    1\leq \kappa_L\leq d.
\end{equation}

Il caso $\kappa_L=1$ corrisponde alla scala della dipendenza asintotica; i
valori $1<\kappa_L<d$ descrivono intermediate tail dependence; il valore
$\kappa_L=d$ coincide con l'ordine dell'indipendenza. A parità delle componenti
lentamente variabili, un valore più piccolo di $\kappa_L$ implica un
decadimento più lento e quindi eventi estremi congiunti relativamente più
probabili.

\subsection{Direzioni diverse dalla diagonale}

Il tail order da solo descrive l'avvicinamento lungo la diagonale. Le tail
order functions estendono l'analisi a direzioni differenti. Se
$\mathbf{w}=(w_1,\ldots,w_d)\in\mathbb{R}_+^d$, la lower tail order function è
definita, quando il limite esiste, da

\begin{equation}
    b(\mathbf{w};C,\kappa_L)
    =\lim_{u\downarrow0}
    \frac{C(uw_1,\ldots,uw_d)}
         {u^{\kappa_L}\ell_L(u)}.
    \label{eq:tail-order-function}
\end{equation}

Questa funzione è omogenea di ordine $\kappa_L$ e descrive come cambia la
probabilità congiunta quando le componenti si avvicinano alla coda con
velocità differenti. Il paper sviluppa proprietà generali del tail order,
studia in particolare copule archimedee multivariate e collega la forma delle
code alla trasformata di Laplace della variabile di mistura. Propone inoltre
una famiglia archimedea basata sulla distribuzione Gamma inversa, capace di
generare configurazioni asimmetriche non riprodotte dalle comuni copule
gaussiane e $t$.

\section{Confronto tra i tre contributi}

La Tabella~\ref{tab:confronto-paper} sintetizza il ruolo dei lavori nel
percorso di studio.

\begin{table}[htbp]
    \centering
    \caption{Confronto tra i tre contributi introduttivi.}
    \label{tab:confronto-paper}
    \begin{tabularx}{\textwidth}{%
        >{\raggedright\arraybackslash}p{0.20\textwidth}%
        >{\raggedright\arraybackslash}p{0.22\textwidth}%
        >{\raggedright\arraybackslash}X}
        \toprule
        Contributo & Strumento centrale & Funzione nel percorso \\
        \midrule
        \textcite{schmidt2005}
        & $\lambda_L$, $\lambda_U$
        & Introduce definizione, calcolo, stima e applicazioni finanziarie del
          coefficiente classico. \\
        \addlinespace
        \textcite{duranteEtAl2015}
        & $\delta_C(t)$ e $q_C(t)$
        & Collega la diagonale della copula alle code e integra il limite
          asintotico con l'andamento a soglie finite. \\
        \addlinespace
        \textcite{huaJoe2011}
        & $\kappa$ e tail order functions
        & Distingue differenti velocità di decadimento e tratta la dipendenza
          intermedia nel caso multivariato. \\
        \bottomrule
    \end{tabularx}
\end{table}

Nel loro insieme, i lavori suggeriscono che la tail dependence non debba
essere trattata come un'unica misura, ma come un insieme di proprietà
complementari. Il coefficiente $\lambda$ offre un riassunto interpretabile del
limite estremo; la tail concentration function mostra quanto rapidamente quel
limite venga raggiunto; il tail order quantifica la forza residua della
dipendenza quando il coefficiente classico è nullo. Questa gerarchia fornisce
una possibile struttura concettuale per la tesi e prepara il successivo studio
degli strumenti di stima e delle applicazioni ai dati finanziari.

\section{Questioni aperte per lo sviluppo della tesi}

La lettura preliminare lascia aperte alcune decisioni che dovranno essere
discusse con il relatore:

\begin{itemize}
    \item se l'obiettivo principale sarà teorico, empirico o misto;
    \item se l'analisi sarà bivariata oppure multivariata;
    \item quali famiglie di copule saranno oggetto di confronto;
    \item se la componente empirica riguarderà rendimenti, perdite, default o
          altre variabili finanziarie;
    \item quali stimatori saranno utilizzati e come verrà scelta la soglia;
    \item se il fulcro sarà il coefficiente classico, il comportamento a scala
          finita o l'intermediate tail dependence.
\end{itemize}

Fino alla definizione della domanda di ricerca, il presente capitolo deve
essere considerato un quadro introduttivo ragionato dei testi assegnati e non
una versione definitiva della rassegna della letteratura.

\printbibliography

\end{document}
